% S6: copied from ../aistats/appendix (see SETUP_REPORT.md for provenance); edit here, not there.
% Trimmed 2026-10-01 (agent C): cut fig:transport-heldout (same data as tab:transport-heldout), tab:moran (final values
% in main fig:programmes d; the without-adversary values kept in the text), tab:kl-summary and fig:kl-seeds (one
% paragraph keeps the negative result), and the Cellina comparison paragraph that repeats main §3.6.
\section{Results on Tissue in Full}
\label{app:full-results}

\subsection{What the Latents Carry and the Response Programmes}
\label{app:response}

\paragraph{Niche information of both latents.} With the $k$-means niches, the response carries $0.41$ to $0.63$ nats of information about the niche beyond its floor and the intrinsic state $0.06$ to $0.11$, by the nearest-neighbour estimator, over three seeds of all four sections; of the floor-corrected total, $85$ to $90\%$ lies in the response. The neural estimator agrees. These estimates use a different cell sample and floor from \cref{tab:headline-full}.

\paragraph{Hallmark labels.} \Cref{fig:hallmarks} shows the enrichment of all $50$ hallmark sets on all $34$ programmes of the twelve fits; $39$ sets are significant on at least one. EMT is significant on $30$ of the $34$ programmes, followed by oestrogen response early ($19$), KRAS signalling up ($17$), TNF$\alpha$/NF-$\kappa$B signalling and oestrogen response late ($15$ each) and coagulation ($14$); $14$ of the $39$ sets are significant on at most two programmes.

\begin{figure*}[t]
\centering
\includegraphics[width=\textwidth]{recomb_fig_hallmarks}
\caption{Hallmark labels of every response programme, by section, seed and programme (within a fit in order of variance; its share of $\vw$'s within-type variance, in $\%$, above). Rows: the sets significant on at least one programme (number of such programmes in parentheses). Cells: AUC of the Mann--Whitney test of the set's absolute loadings against the rest of the expressed panel ($0.5$ is the null); dot, $q \leq 0.05$ (Benjamini--Hochberg within each programme) with AUC above $0.5$. Programmes are identified only up to rotation, so columns are compared within a section by the recurrence of their labels, not as matched axes.}
\label{fig:hallmarks}
\end{figure*}

\paragraph{Localisation.} The first programme leans toward extracellular products (rank-biserial effect $0.08$ to $0.22$) and away from nuclear ones in every seed of the ovarian FFPE, lung and TMA sections; on the fresh-frozen section the extracellular effect is $0.17$, $0.06$ and $0.00$ in the three seeds, fading to null. A gene-specific spill-over of transcripts near the membrane would produce the same lean, and $\kappa$ carries no gene index (\cref{sec:generative}), so the lean does not distinguish a response from spill-over.

\paragraph{Signalling genes.} In the spill-over channel, ligand and receptor genes take a larger share of the between-niche shift than other genes on every section and seed (rank-biserial $0.18$ to $0.28$, $p < 10^{-13}$); in the response channel on the ovarian FFPE, lung and TMA sections ($0.12$ to $0.18$, $p < 10^{-5}$), but not on the fresh-frozen section ($0.01$ to $0.03$, $p = 0.18$ to $0.61$). MintFlow's criterion is therefore met by the spill-over channel alone on every section. The response's lean breaks down at $\kappa = 0.2$ on the fresh-frozen section (\cref{tab:kappa-star}).

\paragraph{True and false tumour axis.} On the primary section the response-predicted shifts follow the tumour bands with a mean $|\tau|$ of $0.78$ to $0.81$ and the false axis with $0.51$ to $0.56$. The contrast holds in tumour cells and macrophages and is reversed in the merged fibroblast class ($0.63$ to $0.71$ against $0.71$ to $0.75$), whose stromal half lies in deep stroma where the false coordinate follows the bands. Shifts predicted from the intrinsic state ($0.59$ to $0.65$ against $0.39$ to $0.47$), from depth ($0.57$ to $0.68$ against $0.36$ to $0.42$) and the observed shift ($0.47$ against $0.33$ to $0.36$) show contrasts of similar size: the response is closest to the bands, but the band ordering is not specific to it.

\paragraph{Marker pairs.} Exclusive double positives are $0.3$ to $2.2\%$ of held-out cells on raw counts and control double positives $1.6$ to $11.8\%$. Decoding alone lowers the exclusive rate on every section and raises the control rate on three (lowers it on the primary section), which sets the sign of this trajectory (\cref{app:sweep}). At $\kappa = 0.1$ the spill-over correction lowers both rates further (exclusive by $0.10$ to $0.50$ points, control by $0.18$ to $1.13$) and their ratio on every section, but the data show no removal specific to exclusive pairs, and the rate falls steadily with $\kappa$ with no plateau, so it does not identify the spill-over fraction.

\subsection{Without the Adversary}
\label{app:moran}

Within-type Moran's I of the posterior means (each coordinate centred on its type's mean, on the model's graph restricted to target cells, averaged with weights equal to the coordinates' within-type variance; permutation null at most $0.005$) shows what the adversary removes. Without it ($\alphaa = 0$, two seeds), I of $\vmu_z$ is $0.09$ to $0.47$; with it, $0.05$ to $0.35$ (\cref{fig:programmes}d), lower by $42$ to $48\%$ on three sections and by $25\%$ on the fresh-frozen one. What remains is within-type variation shared by neighbours that composition and the masked image do not predict (\cref{app:conditional}). Without the adversary the response's posterior mean has almost no within-type variance: summed over coordinates, below $10^{-9}$ in both fresh-frozen fits, $10^{-5}$ and $0.02$ on the ovarian FFPE section and $0.03$ to $0.12$ on the lung and TMA sections, against $3.0$ to $39$ in the final fits. The intrinsic state carries the niche instead.
% source (audit item 14): validation/validation.json morans.mu_w.var_per_dim, summed, of uncontrolledL_s0-s1 and finalL_s0-s2 (target cells, Unassigned mask); devlog 2026-09-29 "Within-type Moran's I table (results; writer)".

\subsection{The Response at the Type Level}
\label{app:kl-maps}

The per-cell divergence of the response from its prior is small, $0.004$ to $0.09$ nats per cell, and differs up to ten-fold between seeds. Within one fit it is spatially organised (Moran's I $0.15$ to $0.85$), but which cells deviate does not reproduce between seeds (per-cell rank correlation $-0.43$ to $0.57$). Under a rule fixed in advance (odds ratio at least $1.5$ or standardised difference at least $0.3$, interval clear of the null, in all three seeds), the top decile differs from the other cells only coarsely: it avoids tumour on the two FFPE sections and lies among fibroblasts (ovarian) and chondrocytes (TMA core); nothing passes on the fresh-frozen section and no continuous feature passes anywhere. The divergence is therefore not a per-cell anomaly score at the operating point, and a planted per-cell response is not detected at this weight either (\cref{app:planted-percell}). The posterior standard deviation of $\vw$ is $0.98$ to $1.01$: the response posterior is its prior in spread as well as in location. The intrinsic state's divergence and standard deviation do reproduce between seeds (rank correlations $0.78$ to $0.93$ and $0.53$ to $0.93$); deeper cells have smaller divergences and wider posteriors (pooled $\rho = -0.77$ to $-0.39$ and $0.21$ to $0.73$), the per-cell scaling of \cref{app:pathb}. The intrinsic state is thus regularised less in low-count cells, and a deep cell's posterior width overstates its uncertainty, one more reason the reported uncertainty comes from the sweep and the seeds rather than from $q$.

\subsection{Relocation in Full}
\label{app:cellina-cf}

\paragraph{Against Cellina.} Cellina's neighbour-rewiring counterfactual \citep{moeed2026cellina} is refitted on the training tiles only, decoded at its posterior mean, and scored by our relocation reads on the panels, niches, ceilings and bootstrap draws of DISCELL's first seed (\cref{tab:cellina-cf}). It has no spill-over channel, so its prediction cannot be split into programme and spill-over parts. The comparison is of one Cellina fit with one DISCELL seed; main \cref{tab:relocation} reads it.

\input{\tablesdir/cellina_cf}

\paragraph{On held-out tiles.} The relocation read is cross-fitted over spatial folds drawn from every tile, so most of its scored cells lie in the model's training tiles. \Cref{tab:transport-heldout} scores the same read on held-out tiles only, with every model quantity still estimated from training tiles and the same panels, ceilings and mask. The rule set before the read was to reconsider the cross-fitted read if the held-out fraction of the ceiling fell outside its interval on any section; it did, in both directions, and the cross-fitted read was kept, because an in-sample advantage would lower the held-out read throughout and the fraction of the ceiling does not move that way. The single-cell read does, on $11$ of $12$ fits. The held-out pool is about three quarters of the cross-fitted one, so fewer panels reach the trusted tier and several all-panel fractions exceed one; on the TMA core no panel is trusted on held-out tiles. On held-out tiles Cellina's fraction of the ceiling stays above DISCELL's range on the ovarian FFPE section and falls below it on the lung section.

\input{\tablesdir/transport_heldout}

\paragraph{Against a plain regression.} On the panels of each seed's relocation read, a ridge regression of the type's expression on its neighbour composition is fitted to the raw counts of the cells that supply the model quantities, and its predicted shift between the two niches is scored by the same fraction of the ceiling. Fitted on $\log(1+x)$ at the median depth, as first specified, it recovers $0.06$ to $0.22$ of the ceiling in the trusted tier; fitted on depth-normalised rates, the scale of the read, it recovers $0.81$ to $0.89$ on the three whole-slide sections (DISCELL $0.65$ to $0.72$) and $0.77$ on the TMA core (DISCELL $0.73$, two seeds with a trusted panel). Over all panels DISCELL is higher on the TMA core and the two FFPE sections, and lower on the FF section ($0.66$ against $0.79$). The single-cell read cannot score the regression against DISCELL's decode of the target cells, so both are scored against the raw target cells: the gap closed is $0.48$ to $0.62$ for the regression and $0.25$ to $0.63$ for DISCELL. Because the observed shifts contain the spill-over, none of these reads can favour a model for separating it (\cref{sec:relocation}); in simulation, scored against the spill-free shift, they can (\cref{app:S-reloc-clean}).
% Source: scripts/logs/recomb_additions_2026-10-01/READOUT.md (3-seed means over finalL_s0-s2; per seed in
%   data/datasets/<ds>/experiments/transport_regression_reference.md). Deviation 2 and 3 of the devlog launch entry.

\subsection{The Latents on Every Section}
\label{app:latents}

\Cref{fig:latents-umap} extends main \cref{fig:latents-main} to all four sections. The type of panels c and d is chosen by a rule: among types with at least $3{,}000$ cells, the one whose neighbourhoods are most mixed. UMAP distorts distances, so the figure illustrates what the probe (\cref{tab:probe}) measures and replaces it nowhere.

\begin{figure}[tp]
\centering
\includegraphics[width=0.8\textwidth]{fig_latents_umap}
\caption{UMAP of the posterior means, one seed of the final configuration; columns: the four sections. \textbf{a},~$\vmu_z$ of $30{,}000$ random cells, coloured by the cell's lineage family. \textbf{b},~$\vmu_w$ of the same cells, coloured by the family dominating its graph neighbours (grey if none). \textbf{c},~$\vmu_z$ of up to $10{,}000$ cells of one type (named above), coloured as in b. \textbf{d},~$\vmu_w$ of the same cells. Families as in \cref{fig:sections}.}
\label{fig:latents-umap}
\end{figure}

\subsection{Training Cost}
\label{app:timing}

\Cref{tab:timing} gives the cost of a fit. DISCELL's time per epoch is pure training; its wall-clock time covers the three final fits including evaluations every five epochs and the early stop, whereas the comparison methods' times are training only, as recorded for the fits compared here. MintFlow's whole-section fits took $14.8$ (lung FFPE) and $15.8$ hours (ovarian FFPE) of wall-clock time, against $3.2$ to $7.4$ minutes for a whole DISCELL fit of the lung section; the other whole-section SIMVI and MintFlow fits could not be run (reasons in \cref{tab:baselines,tab:timing}). The cost of an epoch is linear in the number of cells, and peak memory is the resident count matrix plus the spill-over influx over the edges of a tile, which the tile size bounds (\cref{app:tiles}). The sweep costs eighteen fits per section. Sample size is not analysed beyond the range of sizes in \cref{tab:sections}.

\input{\tablesdir/timing}
